\documentclass{beamer}

\usepackage{packages/packages_mathalea}
\usepackage{fontawesome}

\pgfplotsset{compat=1.18}

\begin{document}

\begin{frame}{Croisement}
Deux villes A et B sont distantes de 40 km. À 10 h, une automobiliste part de A et arrive en B à 10 h 30. Elle fait alors une pause d’une heure puis retourne en A à la même vitesse qu’à l’aller. À 10 h 15, un cycliste part de B vers A à une vitesse de 20 km/h sans s’arrêter. Les mouvements sont supposés uniformes.

\begin{enumerate}
    \item Représenter la situation sur un graphique avec le temps en abscisse et la distance en ordonnée.
    \item Déterminer les équations des trois segments de droites représentant le trajet de l’automobiliste.
    \item Déterminer une équation du segment de droite représentant le trajet du cycliste.
    \item Déterminer graphiquement le nombre de fois où le cycliste et l’automobiliste se croisent.
    \item Déterminer précisément par le calcul le lieu et l’heure de ce(s) croisement(s).
\end{enumerate}
\end{frame}
\end{document}
